\documentclass[final,a4paper,12pt]{article}
\usepackage{notes}

\newcommand{\half}{\tfrac{1}{2}}

\title{Problem Set: Numerical Methods --- Key}
\author{Geophysics 3A}
\date{}

\begin{document}
\maketitle

\begin{enumerate}

%--------------------------------------------------
\item \textbf{Governing equation}

The steady-state heat conduction equation with spatially variable thermal
conductivity and internal heat production is:

\begin{equation}
\nabla \cdot (k \nabla T) + A = 0
\end{equation}

In two dimensions $(x,z)$:

\begin{equation}
\frac{\partial}{\partial x} \left( k \frac{\partial T}{\partial x} \right)
+
\frac{\partial}{\partial z} \left( k \frac{\partial T}{\partial z} \right)
+ A = 0
\end{equation}

%--------------------------------------------------
\item \textbf{Boundary conditions}

\begin{enumerate}

\item \textbf{Top and bottom}

Surface:
\begin{equation}
T(x,0) = T_s
\end{equation}

Basal boundary:
\begin{equation}
T(x,z_b(x)) = T_b(x)
\end{equation}

\item \textbf{Side boundaries (insulating)}

Zero heat flux:
\begin{equation}
\frac{\partial T}{\partial x} = 0
\end{equation}

Finite difference form:
\begin{equation}
T_{i,1} = T_{i,2}, \quad
T_{i,N} = T_{i,N-1}
\end{equation}

\end{enumerate}

%--------------------------------------------------
\item \textbf{Finite difference equation}

Using fluxes across cell faces:

\begin{align}
&\frac{1}{\Delta x^2}
\left[
k_{i+\half,j}(T_{i+1,j} - T_{i,j})
-
k_{i-\half,j}(T_{i,j} - T_{i-1,j})
\right] \\
+ &\frac{1}{\Delta z^2}
\left[
k_{i,j+\half}(T_{i,j+1} - T_{i,j})
-
k_{i,j-\half}(T_{i,j} - T_{i,j-1})
\right]
+ A_{i,j} = 0
\end{align}

Rearranged:

\begin{equation}
T_{i,j} =
\frac{
k_{i+\half,j}T_{i+1,j}/\Delta x^2 +
k_{i-\half,j}T_{i-1,j}/\Delta x^2 +
k_{i,j+\half}T_{i,j+1}/\Delta z^2 +
k_{i,j-\half}T_{i,j-1}/\Delta z^2
+ A_{i,j}
}{
k_{i+\half,j}/\Delta x^2 +
k_{i-\half,j}/\Delta x^2 +
k_{i,j+\half}/\Delta z^2 +
k_{i,j-\half}/\Delta z^2
}
\end{equation}

%--------------------------------------------------
\item \textbf{Heat production at nodes}

\begin{enumerate}

\item General expression:

\begin{equation}
A_{i,j} = \frac{1}{4} \left(
A_{i-\half,j-\half} +
A_{i+\half,j-\half} +
A_{i-\half,j+\half} +
A_{i+\half,j+\half}
\right)
\end{equation}

\item If $A = A(z)$ only:

All four surrounding cells have values depending only on depth:

\begin{equation}
A_{i,j} = \frac{1}{2} (A_{i-\half} + A_{i+\half})
\end{equation}

If the node depth coincides with the midpoint:

\begin{equation}
A_{i,j} \approx A(z_i)
\end{equation}

\end{enumerate}

%--------------------------------------------------
\item \textbf{Thermal conductivity in four directions}

Conductivities at interfaces:

\begin{align}
k_{i+\half,j} &= \frac{1}{2}\left(
k_{i+\half,j-\half} + k_{i+\half,j+\half}
\right) \\
k_{i-\half,j} &= \frac{1}{2}\left(
k_{i-\half,j-\half} + k_{i-\half,j+\half}
\right) \\
k_{i,j+\half} &= \frac{1}{2}\left(
k_{i-\half,j+\half} + k_{i+\half,j+\half}
\right) \\
k_{i,j-\half} &= \frac{1}{2}\left(
k_{i-\half,j-\half} + k_{i+\half,j-\half}
\right)
\end{align}

%--------------------------------------------------
\item \textbf{Porosity-controlled properties}

\begin{enumerate}

\item Thermal conductivity:

\begin{equation}
k = \left( \frac{1-\phi}{k_m} + \frac{\phi}{k_w} \right)^{-1}
\end{equation}

\item Heat production:

\begin{equation}
A = A_0 (1 - \phi)
\end{equation}

\item \textbf{Qualitative behaviour}

\begin{itemize}
\item Porosity decreases with depth
\item Conductivity increases with depth (less fluid, more solid matrix)
\item Heat production increases with depth (more radioactive solid material)
\end{itemize}

\end{enumerate}

%--------------------------------------------------
\item \textbf{Surface heat flow}

Heat flow given by Fourier's law:

\begin{equation}
q = -k \frac{\partial T}{\partial z}
\end{equation}

Finite difference approximation:

\begin{equation}
q_j \approx -k_{1,j} \frac{T_{2,j} - T_{1,j}}{\Delta z}
\end{equation}

%--------------------------------------------------
\item \textbf{Conceptual question}

Although $k$ and $A$ vary only with depth, the geometry of the system
(sediment thickness) varies laterally.

This creates lateral temperature gradients, which cause heat to flow
horizontally as well as vertically.

Therefore, the system is two-dimensional because:
\begin{itemize}
\item boundary geometry varies in $x$
\item isotherms deform around basement topography
\item heat is redistributed laterally
\end{itemize}

\textbf{Key point:} Geometry, not material properties, drives the 2-D behaviour.

%--------------------------------------------------
\item \textbf{Sketching exercise}

\textbf{1-D case:}
\begin{itemize}
\item Flat sediment layer
\item Horizontal isotherms
\item Vertical heat flow
\end{itemize}

\textbf{2-D case:}
\begin{itemize}
\item Variable sediment thickness
\item Isotherms bend upward over highs
\item Heat flow diverges over highs and converges in basins
\end{itemize}

%--------------------------------------------------
\item \textbf{Model comparison}

\begin{itemize}
\item 1-D model provides smooth, averaged behaviour
\item 2-D model captures local variations due to bathymetry
\item 2-D model fits short-wavelength variability better
\item Both models may miss some systematic trends
\end{itemize}

%--------------------------------------------------
\item \textbf{3-D effects}

\begin{itemize}
\item 2-D ridge: heat deflection occurs in one horizontal direction
\item Seamount (3-D): heat spreads radially
\item 3-D spreading reduces peak anomalies
\end{itemize}

Including 3-D effects can improve the model but will not fully remove misfit.

%--------------------------------------------------
\item \textbf{Improving the model}

Possible improvements:

\begin{itemize}
\item Include hydrothermal circulation
\item Allow lateral variation in material properties
\item Improve boundary conditions (e.g.\ basal temperature structure)
\item Use fully 3-D geometry
\end{itemize}

\end{enumerate}

\end{document}